\section*{Chapter \ref{ch:b2:continuous-reactors} --- Continuous Reactors and Industrial Processes}

\begin{solution}{exo:b2:continuous-reactors:1}
$\tau = V/Q_v = 2000/0.50 = \qty{4.0e3}{s}$, that is \qty{1.1}{h}.
\end{solution}

\begin{solution}{exo:b2:continuous-reactors:2}
$k\tau = 1.2 \times 10^{-3} \times 1800 = 2.16$. Tank: $X = 2.16/3.16 = 0.68$;
tube: $X = 1 - \mathrm e^{-2.16} = 0.88$.
\end{solution}

\begin{solution}{exo:b2:continuous-reactors:3}
Tube: $\tau = \ln 100/0.020 = \qty{230}{s}$. Tank: $X/(1 - X) = 99 = k\tau$,
$\tau = \qty{4950}{s}$, more than twenty times longer: at high \omterm{def:b2:continuous-reactors:conversion}{conversion} the
stirred tank pays dearly for working at the outlet concentration.
\end{solution}

\begin{solution}{exo:b2:continuous-reactors:4}
Per litre: $120 \times 0.50 = \qty{60}{kJ}$ released, heating
\qty{4.2}{kJ/K}: $\Delta T_{\text{ad}} = \qty{14}{K}$. A failure of the cooling
would warm the mixture noticeably but not dangerously, unless the solvent is
near its boiling point or a second reaction starts in that range.
\end{solution}

\begin{solution}{exo:b2:continuous-reactors:5}
Tank: $Da = X/(1 - X)^2 = 0.8/0.04 = 20$, $\tau = 20/0.010 = \qty{2000}{s}$.
Tube: $Da = X/(1 - X) = 4$, $\tau = \qty{400}{s}$. Ratio 5, against $4/\ln 5 =
2.5$ for order 1: the higher the order, the more the stirred tank suffers
from its low concentration.
\end{solution}

\begin{solution}{exo:b2:continuous-reactors:6}
$k\tau = 5.0$. Three tanks: $X = 1 - (1 + 5/3)^{-3} = 0.947$. One tank: $5/6 =
0.833$. Tube: $1 - \mathrm e^{-5} = 0.993$.
\end{solution}

\begin{solution}{exo:b2:continuous-reactors:7}
$X = 0.70$ means $k\tau = 0.70/0.30 = 2.33$. Doubling the flow halves $\tau$:
$k\tau = 1.17$, $X = 0.54$. To recover 70\,\% the \omterm{def:b2:continuous-reactors:residence-time}{residence time}, hence the
volume, must be doubled.
\end{solution}

\begin{solution}{exo:b2:continuous-reactors:8}
Released: $80 \times 2.0 \times 1.0 \times 0.90 = \qty{144}{kW}$. Carried by the
outlet stream, heated from \qty{330}{K} to \qty{350}{K}: $4.0 \times 1.0 \times 20 =
\qty{80}{kW}$. The jacket must remove the difference, \qty{64}{kW}.
\end{solution}

\begin{solution}{exo:b2:continuous-reactors:9}
Linear dimensions $\times 10$: volume $\times 1000$, wall area $\times 100$, area
per unit volume $\div 10$. The heat released grows with the volume, the heat
removed through the wall with its area: the large reactor removes, per litre,
ten times less heat at the same temperature difference.
\end{solution}

\begin{solution}{exo:b2:continuous-reactors:10}
Tube: $\tau = \ln(0.05/0.10)/(0.05 - 0.10) = \qty{13.9}{min}$, $c_B/c_{A,0} =
\frac{0.10}{-0.05}(\mathrm e^{-1.386} - \mathrm e^{-0.693}) = 0.50$. Tank: $\tau =
1/\sqrt{0.005} = \qty{14.1}{min}$, $c_B/c_{A,0} = 1.414/(2.414 \times 1.707) =
0.34$. In the tank, fresh A and finished B are mixed in the same volume: some
B always stays long enough to be converted to C while some A has barely
arrived; in the tube every slice has the same history.
\end{solution}

\begin{solution}{exo:b2:continuous-reactors:11}
Raising the coolant (and feed) temperature moves the removal line to the
right. The cold intersection rises slowly until the line becomes tangent to
the lower bend of the S; beyond, the cold state disappears and the tank jumps
to the hot branch (ignition). Lowering the temperature again, the tank stays
on the hot branch until the line becomes tangent to the upper bend, and only
then falls back (extinction), at a lower temperature than ignition: a
hysteresis loop.
\end{solution}

\begin{solution}{exo:b2:continuous-reactors:12}
Batch: 75\,\% of \qty{2.0}{mol/L} remain, $1.5 \times 100 = \qty{150}{kJ/L}$,
$\Delta T = 150/4.0 = \qty{37.5}{K}$ above the working temperature, enough to
reach a boiling point or a decomposition. Semi-batch: at most
\qty{0.10}{mol/L} unreacted, $\Delta T = 10/4.0 = \qty{2.5}{K}$. Slow addition
caps the energy stored in the reactor.
\end{solution}

\begin{solution}{pb:b2:continuous-reactors:1}
\textbf{1.} $r = k[\text{ester}][\ce{OH-}]$; equal feeds and a 1 : 1
stoichiometry keep the two concentrations equal, $r = kc^2$.
\textbf{2.} $1/c - 1/c_0 = kt$ with $c = 0.10c_0$: $kc_0t = 9$, $t = 9/(0.11
\times 0.10) = \qty{818}{s}$, about \qty{14}{min}.
\textbf{3.} $t_{1/2} = 1/(kc_0) = \qty{91}{s}$; for order 2 the rate falls with
the square of the concentration, so the time to halve it depends on where one
starts.
\textbf{4.} $Da = kc_0\tau = 0.011\,\tau$ ($\tau$ in seconds).
\textbf{5.} $Q_vc_0 - Q_vc - kc^2V = 0$, that is $c_0 - c = k\tau c^2$.
\textbf{6.} With $c = c_0(1 - X)$: $c_0X = k\tau c_0^2(1 - X)^2$.
\textbf{7.} $Da = 0.90/0.10^2 = 90$.
\textbf{8.} $\tau = 90/0.011 = \qty{8.2e3}{s}$, \qty{2.3}{h}.
\textbf{9.} $V = 0.50 \times 8182 = \qty{4.1e3}{L}$, \qty{4.1}{m^3}.
\textbf{10.} The outlet concentration, $0.010\,\unit{mol/L}$: the rate is
everywhere a hundred times lower than at the inlet.
\textbf{11.} $Da = X/(1 - X) = 9$, $\tau = \qty{818}{s}$, $V = \qty{409}{L}$.
\textbf{12.} Ten times less volume for the tube.
\textbf{13.} $c_{j-1} - c_j = k\tau_1c_j^2$, with $\tau_1$ the \omterm{def:b2:continuous-reactors:residence-time}{residence time} of
one tank.
\textbf{14.} $\tau_1 = 850/(3 \times 0.50) = \qty{567}{s}$, $k\tau_1 =
\qty{62.4}{L/mol}$. Solving $62.4c_j^2 + c_j - c_{j-1} = 0$: $c_1 = 0.0328$,
$c_2 = 0.0163$, $c_3 = \qty{0.0100}{mol/L}$: 90\,\% \omterm{def:b2:continuous-reactors:conversion}{conversion}.
\textbf{15.} The tube, or the cascade (ten and five times smaller than one
tank). A single tank may still win for its uniform temperature, easy
cleaning and control, and because here volume is cheap.
\textbf{16.} $\Delta T_{\text{ad}} = 55\,000 \times 100/(997 \times 4180) =
\qty{1.3}{K}$ (\qty{1.2}{K} at 90\,\%).
\textbf{17.} $55 \times 0.50 \times 0.10 \times 0.90 = \qty{2.5}{kW}$.
\textbf{18.} No: the outlet stream carries the heat away with a rise of about
a kelvin; no runaway is possible.
\textbf{19.} Twenty times more concentrated: the adiabatic rise becomes
\qty{26}{K} and the heat \qty{50}{kW}, worth a cooling coil; the rate $kc_0$
being twenty times higher, the tank for 90\,\% would be twenty times smaller
(\qty{0.20}{m^3}).
\textbf{20.} $0.50 \times 0.10 \times 0.90 = \qty{0.045}{mol/s}$, $\times 86\,400
\times 82.03 = \qty{319}{kg}$ of sodium ethanoate per day.
\textbf{21.} The rate constant grows with temperature (Arrhenius), so the
volume shrinks; the cost is the heating of the feed and, for other
reactions, a lower \omterm{def:b2:continuous-reactors:selectivity}{selectivity} and more vapour.
\textbf{22.} The single stirred tank must hold $\boldsymbol{V \approx
\qty{4.1}{m^3}}$.
\end{solution}
